\documentclass[11pt,a4paper]{article}
\usepackage[utf8]{inputenc}
\usepackage[T1]{fontenc}
\usepackage{amsmath,amssymb}
\usepackage{bm}
\usepackage{geometry}
\usepackage{hyperref}
\usepackage{amsthm}
\usepackage{booktabs}
\usepackage{microtype}
\newtheorem{veta}{Veta}
\newtheorem{definicia}{Definícia}
\newtheorem{poznamka}{Poznámka}
\geometry{margin=2.7cm}
\hypersetup{hidelinks}
\renewcommand{\abstractname}{Abstrakt}
\renewcommand{\proofname}{Dôkaz}
\renewcommand{\tablename}{Tabuľka}
\renewcommand{\refname}{Referencie}

\title{Rozklad podľa centra momentov (MVDC):\\
Centrálne Bernoulliho čísla a vysoko presné korekcie Eulerových súčinov}
\author{Ing.~Robert~Polák\\\small e--mail: robopol@gmail.com}
\date{11.~augusta~2026}

\begin{document}
\maketitle

\begin{abstract}
Metóda rozkladu podľa centra momentov (MVDC) rozdeľuje konečný súčin na
člen určený zvoleným centrom a na presné rezíduum. Pri hladkých
postupnostiach faktorov je Taylorov rozvoj logaritmu okolo stredného
indexu riadený centrovanými mocninami indexov. Tieto konečné koeficienty
tvoria prirodzenú \(C\)--Bernoulliho rodinu a majú presné vyjadrenie
pomocou centrálnych Bernoulliho polynómov aj kompaktnú generujúcu
funkciu. Druhý, odlišný momentový rozvoj vyplýva z Taylorovho radu
\(\log(1+x)\), keď multiplikatívne centrujeme samotné faktory, nie ich
logaritmy. Dôsledným oddelením týchto dvoch konštrukcií odstraňujeme
nejednoznačnosť skoršej verzie a zachovávame centrálnu Bernoulliho
štruktúru MVDC. Aplikácia na faktoriál dáva obvyklé presné Bernoulliho
koeficienty okolo analytického centra MVDC. Pri Eulerových súčinoch sa
konečný prvočíselný blok vyhodnotí presne pomocou jeho logaritmického
priemeru a modelovať alebo ohraničiť treba iba zostávajúci nekonečný
chvost.
\end{abstract}

\section{Presný rozklad a úloha centra}

Nech
\[
  P_n=\prod_{i=1}^{n}a_i,\qquad a_i>0,\qquad
  \ell_i:=\log a_i .
\]
Pre ľubovoľné reálne centrum \(c\), kde \(k=e^c\), definujme
\[
  H_c:=k^n=e^{nc},\qquad
  R_c(n):=\sum_{i=1}^{n}(\ell_i-c).
\]
Základnou identitou MVDC je
\begin{equation}\label{eq:presne-centrum}
  \boxed{\;P_n=H_c\,e^{R_c(n)},\qquad
  \log P_n=nc+R_c(n).\;}
\end{equation}
V \eqref{eq:presne-centrum} sa nepoužíva nijaký asymptotický predpoklad.
Centrum možno zvoliť analyticky tak, aby \(H_c\) absorbovalo známe
dominantné členy, alebo numericky kvôli stabilite.

Existuje jeden dôležitý hraničný prípad. Ak
\[
  c=\mu_{\log}:=\frac1n\sum_{i=1}^{n}\ell_i,
\]
potom \(R_{\mu_{\log}}(n)=0\) a \(H_{\mu_{\log}}=P_n\) identicky.
Ide o platné presné vyhodnotenie v logaritmickej oblasti, ale nezostáva
pri ňom žiadne nenulové rezíduum na aproximáciu. Netriviálna asymptotická
aplikácia preto používa centrum určené nezávisle od presnej hodnoty
konečného súčinu, napríklad integrálom, výpočtom na krajoch intervalu
alebo známou vedúcou asymptotikou.

\section{Korektný rozvoj pomocou faktorových momentov}

Vyššie momenty vstupujú cez samotné faktory. Píšme
\[
  a_i=k(1+x_i),\qquad x_i=\frac{a_i}{k}-1,
  \qquad X_r(k):=\sum_{i=1}^{n}x_i^r .
\]
Ak \(\rho:=\max_i|x_i|<1\), sčítaním Taylorových radov pre
\(\log(1+x_i)\) dostaneme
\begin{equation}\label{eq:faktorove-momenty}
  \boxed{\;
  \log\frac{P_n}{k^n}
  =\sum_{r=1}^{\infty}\frac{(-1)^{r-1}}{r}\,X_r(k).\;}
\end{equation}
Po odrezaní pri \(r=M\) platí explicitný odhad zvyšku
\[
  |\mathcal R_M|
  \le
  \sum_{i=1}^{n}
  \frac{|x_i|^{M+1}}{(M+1)(1-|x_i|)}.
\]
Pre centrum dané aritmetickým priemerom
\(k=\bar a=n^{-1}\sum_i a_i\) platí \(X_1(k)=0\), takže rad začína pri
\(r=2\). Toto je presný prípad, v ktorom zrušenie prvého momentu
vyplýva z radu \(\log(1+x)\). Centrálne logaritmické momenty
\(\sum(\ell_i-\mu_{\log})^r\) sú užitočnými charakteristikami dát,
nesmú sa však zamieňať s faktorovými momentmi \(X_r(k)\) v
\eqref{eq:faktorove-momenty}.

\section{Centrálne Bernoulliho a Nörlundove polynómy}

Obyčajné Bernoulliho polynómy sú definované vzťahom
\[
  \frac{t e^{xt}}{e^t-1}
  =\sum_{r=0}^{\infty}B_r(x)\frac{t^r}{r!}.
\]
Generalizované Bernoulliho polynómy rádu \(\alpha\), spájané s
Nörlundom, sú definované vzťahom
\begin{equation}\label{eq:norlund-sk}
  \left(\frac{t}{e^t-1}\right)^\alpha e^{xt}
  =\sum_{r=0}^{\infty}B_r^{(\alpha)}(x)\frac{t^r}{r!}.
\end{equation}
Jedna zo štandardných centrálnych Bernoulliho rodín je
\begin{equation}\label{eq:centralne-bernoulli-sk}
  B_r^c(x):=2^rB_r\!\left(\frac{x+1}{2}\right),
  \qquad
  \frac{t}{\sinh t}e^{xt}
  =\sum_{r=0}^{\infty}B_r^c(x)\frac{t^r}{r!}.
\end{equation}
Platí \(B_r^c(-x)=(-1)^rB_r^c(x)\). Všeobecnejšie,
\[
  \left(\frac{t}{\sinh t}\right)^\alpha e^{xt}
  =\sum_{r=0}^{\infty}B_{r,\alpha}^c(x)\frac{t^r}{r!},
  \qquad
  B_{r,\alpha}^c(x)
  =2^rB_r^{(\alpha)}\!\left(\frac{x+\alpha}{2}\right).
\]
Centrálne Bernoulliho objekty a ich vzťah k Nörlundovým polynómom sú
teda štandardné; zostáva presne určiť, ako vstupujú do konkrétneho
rozvoja MVDC.

\section{Konečné C--Bernoulliho koeficienty}

Položme \(m_n=(n+1)/2\). Konečné koeficienty centrovaných indexov
použité v tejto práci definujeme takto.
\begin{definicia}
\begin{equation}\label{eq:C-definicia-sk}
  \boxed{\;
  C_r(n):=\frac1{r!}\sum_{i=1}^{n}(i-m_n)^r,\qquad r\ge0.\;}
\end{equation}
\end{definicia}
Definícia je čisto algebraická a závisí iba od geometrie \(n\)
rovnomerne vzdialených indexov. Faulhaberova identita dáva
\begin{equation}\label{eq:C-bernoulli-sk}
  C_r(n)
  =
  \frac{
    B_{r+1}\!\left(\frac{n+1}{2}\right)
    -B_{r+1}\!\left(\frac{1-n}{2}\right)}
  {(r+1)r!}
  =
  \frac{B_{r+1}^c(n)-B_{r+1}^c(-n)}
  {2^{r+1}(r+1)r!}.
\end{equation}
Generujúca funkcia je regulárna v počiatku a má kompaktný tvar
\begin{equation}\label{eq:C-generujuca-sk}
  \boxed{\;
  \sum_{r=0}^{\infty}C_r(n)t^r
  =\sum_{i=1}^{n}e^{(i-m_n)t}
  =\frac{\sinh(nt/2)}{\sinh(t/2)}.\;}
\end{equation}
Zo symetrie teraz vyplýva správna parita:
\[
  C_{2j+1}(n)=0.
\]
Prvé nenulové koeficienty sú
\[
  C_0(n)=n,\qquad
  C_2(n)=\frac{n(n^2-1)}{24},\qquad
  C_4(n)=\frac{n(n^2-1)(3n^2-7)}{5760}.
\]

\begin{veta}\label{veta:stredovy-rozvoj}
Nech \(f\in C^{2p+2}[1,n]\). Potom
\begin{equation}\label{eq:stredovy-rozvoj}
  \sum_{i=1}^{n}f(i)
  =
  \sum_{j=0}^{p}f^{(2j)}(m_n)C_{2j}(n)
  +\mathcal T_{2p+2},
\end{equation}
kde
\[
  |\mathcal T_{2p+2}|
  \le
  \left\|f^{(2p+2)}\right\|_{\infty}
  C_{2p+2}(n).
\]
\end{veta}
\begin{proof}
Na každé \(f(i)\) použijeme Taylorovu vetu v bode \(m_n\), konečné
rozvoje sčítame a použijeme \eqref{eq:C-definicia-sk}. Všetky nepárne
mocniny sa vďaka symetrii zrušia. Sčítaním Lagrangeových odhadov zvyšku
dostaneme uvedenú nerovnosť.
\end{proof}

Pre \(f(i)=\log a_i\) je Veta~\ref{veta:stredovy-rozvoj} rigoróznym
centrálnym Bernoulliho mostom v MVDC. Koeficienty \(C_r(n)\) kódujú
centrované indexy, zatiaľ čo derivácie hladkej interpolácie \(f(x)\)
kódujú konkrétne faktory. Tie isté \(C_r(n)\) preto možno použiť pre
rôzne hladké súčiny bez toho, aby sme ich stotožnili s momentmi hodnôt
\(f(i)\), ktoré závisia od dát.

\section{Eulerov--Maclaurinov tvar a faktoriál}

Pre \(g_c(x)=f(x)-c\) má Eulerova--Maclaurinova formula na indexoch
\(1,\ldots,n\) tvar
\begin{align}
  \sum_{i=1}^{n}g_c(i)
  ={}&
  \int_1^n g_c(x)\,dx
  +\frac{g_c(1)+g_c(n)}2 \notag\\
  &+\sum_{j=1}^{p}
  \frac{B_{2j}}{(2j)!}
  \left(g_c^{(2j-1)}(n)-g_c^{(2j-1)}(1)\right)
  +\mathcal E_p . \label{eq:EM-sk}
\end{align}
Štandardný zvyšok \(\mathcal E_p\) je integrál obsahujúci periodický
Bernoulliho polynóm a \(g_c^{(2p)}\). Centrum môže v konkrétnej aplikácii
absorbovať časť integrálnych alebo okrajových príspevkov, ale tieto
zrušenia treba pre danú aplikáciu osobitne overiť.

Pre faktoriál zvoľme analytické centrum
\[
  c_n=\log n-1+\frac{\log(2\pi n)}{2n},\qquad
  H_n=e^{nc_n}=\sqrt{2\pi n}\left(\frac ne\right)^n .
\]
Eulerova--Maclaurinova formula aplikovaná na \(f(x)=\log x\), so
štandardným určením limitnej konštanty, dáva
\begin{equation}\label{eq:faktorial-sk}
  \log\frac{n!}{H_n}
  \sim
  \sum_{j=1}^{\infty}
  \frac{B_{2j}}{2j(2j-1)n^{2j-1}}.
\end{equation}
Teda
\[
  \log\frac{n!}{H_n}
  =
  \frac1{12n}-\frac1{360n^3}
  +\frac1{1260n^5}-\frac1{1680n^7}
  +O(n^{-9}).
\]
Sú to presne analytické koeficienty implementované v súbore
\texttt{mvdc\_factorial\_analytic.py}. Pre toto konkrétne centrum sú
usporiadanie MVDC a klasický logaritmický Stirlingov rozvoj tou istou
uzavretou formulou zapísanou ako \(H_n\) krát malá korekcia. Toto
pozorovanie neporovnáva iné centrá MVDC ani fitované kaskády so
Stirlingovými rozvojmi.

\begin{table}[ht]
\centering
\begin{tabular}{rc}
\toprule
\(n\) & relatívna chyba po člene \(n^{-7}\) \\
\midrule
10  & \(8.23\times10^{-13}\) \\
20  & \(1.63\times10^{-15}\) \\
50  & \(4.31\times10^{-19}\) \\
100 & \(8.42\times10^{-22}\) \\
\bottomrule
\end{tabular}
\caption{Vysoko presná kontrola \eqref{eq:faktorial-sk} s 80-cifernou aritmetikou.}
\end{table}

\section{Aplikácia na orezané Eulerove súčiny}

Nech
\[
  \zeta_N(s)=\prod_{p\le N}(1-p^{-s})^{-1},
  \qquad \sigma:=\Re s>1.
\]
Pre \(M>N\) definujme
\[
  f_s(p):=-\log(1-p^{-s}),\qquad
  \mu_{N,M}(s):=\frac1n\sum_{N<p\le M}f_s(p),
  \qquad n=\pi(M)-\pi(N).
\]
Konečný blok je presne
\[
  B_{N,M}(s):=\sum_{N<p\le M}f_s(p)=n\mu_{N,M}(s).
\]
Presný rozklad chvosta je preto
\begin{equation}\label{eq:zeta-chvost-sk}
  \log\frac{\zeta(s)}{\zeta_N(s)}
  =n\mu_{N,M}(s)+R_M(s),
  \qquad
  R_M(s)=\sum_{p>M}-\log(1-p^{-s}).
\end{equation}
Pre vynechaný chvost platí elementárny rigorózny odhad
\begin{equation}\label{eq:odhad-chvosta-sk}
  |R_M(s)|
  \le
  \frac{M^{1-\sigma}}
  {(\sigma-1)(1-2^{-\sigma})}.
\end{equation}
Skutočne,
\(-\log(1-p^{-\sigma})\le p^{-\sigma}/(1-2^{-\sigma})\)
a súčet cez prvočísla zhora ohraničíme zodpovedajúcim integrálom cez
celé čísla.

Na numerické zrýchlenie možno \(R_M(2)\) modelovať integrálom
\[
  I_M:=\int_M^\infty\frac{x^{-2}}{\log x}\,dx.
\]
Ide o model hustoty prvočísiel, nie o rigorózny odhad
\eqref{eq:odhad-chvosta-sk}. Pre \(M=10N\) dáva priamy vysoko presný
výpočet:

\begin{table}[ht]
\centering
\begin{tabular}{rrrrr}
\toprule
\(N\) & \(n\) & presný chvost \(\Delta\) & \(B_{N,M}+I_M\) & chyba \\
\midrule
1000  & 1061 & \(1.26989814756\times10^{-4}\)
              & \(1.27049267685\times10^{-4}\) & \(5.95\times10^{-8}\) \\
5000  & 4464 & \(2.10858622405\times10^{-5}\)
              & \(2.10884676662\times10^{-5}\) & \(2.61\times10^{-9}\) \\
10000 & 8363 & \(9.81570113247\times10^{-6}\)
              & \(9.81712884998\times10^{-6}\) & \(1.43\times10^{-9}\) \\
\bottomrule
\end{tabular}
\caption{Presný konečný blok MVDC plus integrálny model zostávajúceho zeta chvosta.}
\end{table}

Numerická presnosť v tabuľke je skutočná, ale pochádza z presného
konečného bloku a z integrálneho modelu chvosta. V skoršom experimente
mal dodatočný člen z centrálnych logaritmických momentov pri \(N=1000\)
hodnotu \(-1.56\times10^{-14}\), čo je hlboko pod zobrazenou chybou
\(5.95\times10^{-8}\). Takéto momenty možno naďalej skúmať ako
diagnostické veličiny, nie však ako dodatočný člen v
\eqref{eq:zeta-chvost-sk}.

\subsection{Rozšírenie na Dirichletove \texorpdfstring{\(L\)}{L}-funkcie}

Pre primitívny Dirichletov charakter \(\chi\) modulo \(q\) definujme
\[
  f_{\chi,s}(p):=-\log(1-\chi(p)p^{-s}),\qquad p\nmid q.
\]
Pri vetve logaritmu zhodnej s absolútne konvergentným Eulerovým súčinom
platí ten istý presný rozklad:
\begin{equation}\label{eq:L-chvost-sk}
  \log\frac{L(s,\chi)}{L_N(s,\chi)}
  =
  n\mu_{N,M}(\chi,s)+R_{M,\chi}(s),
\end{equation}
kde \(n\) počíta prvočísla \(N<p\le M\) s \(p\nmid q\). Keďže
\(|\chi(p)|\le1\), odhad \eqref{eq:odhad-chvosta-sk} ohraničuje aj
\(|R_{M,\chi}(s)|\).

Pre netriviálny reálny charakter modulo \(3\),
\[
  \chi_3(n)=
  \begin{cases}
    1,&n\equiv1\pmod3,\\
   -1,&n\equiv2\pmod3,\\
    0,&3\mid n,
  \end{cases}
\]
možno referenčnú hodnotu nezávisle vypočítať zo vzťahu
\[
  L(2,\chi_3)
  =\frac1{9}
   \left[\zeta\!\left(2,\frac13\right)
        -\zeta\!\left(2,\frac23\right)\right].
\]
Pre \(N=10^3\) a \(M=10^6\) dáva 60-ciferná aritmetika
\[
\begin{array}{lr}
\text{presný chýbajúci chvost} & -5.61728190475735\times10^{-8},\\
\text{presný konečný blok }(N,M] & -5.61013233973667\times10^{-8},\\
\text{vynechaný chvost }p>M & -7.14956502068784\times10^{-11}.
\end{array}
\]
Predtým pozorovaná zhoda na úrovni \(10^{-10}\) teda zostáva zachovaná,
pričom jej zdrojom je správne identifikovaný explicitne vypočítaný
prvočíselný blok, nie pridaný rad centrálnych momentov.

\section{Záver}

Opravená analytická štruktúra MVDC má štyri odlišné časti:
\par\smallskip
\begin{itemize}
  \item presný rozklad podľa centra \eqref{eq:presne-centrum}, platný pre
        každý kladný konečný súčin;
  \item konvergentný rozvoj pomocou faktorových momentov
        \eqref{eq:faktorove-momenty} s explicitným zvyškom;
  \item konečné \(C\)--Bernoulliho koeficienty
        \eqref{eq:C-definicia-sk}, ktoré kódujú centrované súčty indexov
        a vstupujú do hladkých stredových rozvojov cez
        Vetu~\ref{veta:stredovy-rozvoj};
  \item analytické alebo numerické modely rezídua špecifické pre danú
        aplikáciu, napríklad \eqref{eq:faktorial-sk} a
        \eqref{eq:zeta-chvost-sk}.
\end{itemize}
Centrálne Bernoulliho a Nörlundove polynómy tak zostávajú skutočnou
súčasťou rámca MVDC. Nemožno však zamieňať koeficienty centrovaných
indexov, momenty odchýlok faktorov a centrálne momenty logaritmických
dát. Po ich dôslednom oddelení zostávajú platné faktoriálové aj
Eulerovské výpočty zachované a zdroje ich chýb možno uviesť presne.

\begin{thebibliography}{9}
\bibitem{DLMF-EM-sk}
NIST Digital Library of Mathematical Functions,
\emph{Euler--Maclaurin Formula}, Sec.~2.10(i),
\url{https://dlmf.nist.gov/2.10}.

\bibitem{DLMF-Bernoulli-sk}
NIST Digital Library of Mathematical Functions,
\emph{Bernoulli and Euler Polynomials: Sums of Powers}, Sec.~24.4(iii),
\url{https://dlmf.nist.gov/24.4}.

\bibitem{DLMF-Generalized-sk}
NIST Digital Library of Mathematical Functions,
\emph{Generalized Bernoulli and Nörlund Polynomials}, Sec.~24.16,
\url{https://dlmf.nist.gov/24.16}.

\bibitem{Tempesta-sk}
P.~Tempesta,
\emph{On Appell sequences of polynomials of Bernoulli and Euler type},
\emph{Journal of Mathematical Analysis and Applications}
\textbf{341} (2008), 1295--1310.

\bibitem{Norlund-sk}
N.~E.~Nörlund,
\emph{Vorlesungen über Differenzenrechnung},
Springer, Berlin, 1924.
\end{thebibliography}

\end{document}
